\scanpage{87}
se factorise en
\[
\tag{121}
\diagram{\begin{tikzcd}[ampersand replacement=\&, column sep=large]
\underbrace{(S_C\Pi M_G,\Gamma^!_{G,G'}\deleted{\unsure{\prime}})}_{\deleted{\Pi M^*_{G,G'}}} \arrow[r] \& \Pi M_{G'},
\end{tikzcd}}
\]
\deleted{que}\inserted{et} plus généralement, si $\Gamma^!_{G,G'}\subset\Gamma'\subset\Gamma_{G,G'}$ est un sous-groupe de $\Gamma=\operatorname{Aut}G$ qui stabilise $C$ et contient $\Gamma^!_{G,G'}$, on trouve
\[
\tag{122}
\diagram{\begin{tikzcd}[ampersand replacement=\&, column sep=large,row sep=large]
(S_{\unsure{S|C}}\Pi M_G,\Gamma') \arrow[r] \arrow[d,equal] \& (\Pi M_{G'},\Gamma'/\Gamma^!_{G,G'})\\
(\Pi M_{G,G'},\Gamma'/\Gamma^!_{G,G'})
\end{tikzcd}}
\]

Ici la relation de transitivité un peu compliquée \inserted{(105)} s'exprime, plus simplement, par le fait que les \inserted{groupoïdes} $S_A\Pi M_G$ dépendent fonctoriellement des \inserted{MD-}graphes $G$, pour les morphismes de générisation.

On va axiomatiser un peu. \deleted{plus.} Les MD-graphes \inserted{et morphismes de générisation} forment une \uline{catégorie}, $(\deleted{\unsure{\mathcal G}}\mathfrak g)$, où tous les morphismes sont des épimorphismes. Cette catégorie satisfait
\margin{La catégorie $\mathcal G$ joue le rôle d'un ens. de \unsure{complexités|« types »} géométriques \unsure{ou des cas}\ldots}
\margin{86}
